一份 Bayesian 分析交出的后验是一张完整的参数分布，可它的归一化常数是个一百维积分。被积函数在绝大多数区域贴近零，只在几个窄峰上极高；均匀网格每维取十个点就是 \(10^{100}\) 次求值，这个数字比可观测宇宙的原子数还大。同样这个积分，取 \(10^{5}\) 个随机样本，标准误约 \(0.3/\sqrt{10^{5}}\approx0.001\)，而且这个精度与维数无关，只随样本量走。整个单元从这处落差开始。

落差的两侧不对称：目标量的表达式是完整的，通往数字的路却被维数堵死。类似的堵点在概率模型里到处都是——后验的归一化常数、能量模型的配分函数、隐变量的边缘化、状态空间的滤波递推。绕过去的路只有两条。

\begin{center}
\begin{tikzpicture}[>=stealth,every node/.style={font=\small}]
\node[draw,align=center] (a) at (0,0) {目标量写得出\\积分算不出};
\node[draw,align=center] (b) at (4.2,1.7) {采样\\用样本的经验分布逼近};
\node[draw,align=center] (c) at (4.2,-1.7) {优化\\用一族简单分布逼近};
\node[align=center] (d) at (9.4,1.7) {代价：Monte Carlo 误差、\\自相关、跨不过模态};
\node[align=center] (e) at (9.4,-1.7) {代价：族外的形状表达不出、\\反向 KL 压窄方差};
\draw[->] (a) -- (b);
\draw[->] (a) -- (c);
\draw[->] (b) -- (d);
\draw[->] (c) -- (e);
\end{tikzpicture}
\end{center}

\emph{两条路都不是免费的，选哪条取决于你更付得起哪种代价。}

采样这条路上，能独立抽样时用样本平均直接估期望，误差有 \(n^{-1/2}\) 的明码标价；不能独立抽样时——后验通常如此，密度只知道差一个常数——就构造一条以目标分布为平稳分布的 Markov 链。要留心的是链的合法性来自遍历性而不是样本独立：时间平均之所以收敛到空间平均，是因为链按目标分布的比例访问了整个空间。相邻样本当然相关，那只影响精度，不影响相合性；而收敛快慢由转移算子的谱间隙决定，burn-in 长度与自相关时间是同一件事的两种说法。链没跨过模态时，你手里那条平稳得毫无破绽的轨道其实被困在一个峰里，识破它只能靠多链、\(\widehat R\) 与有效样本量，靠不了迭代步数。

优化这条路把积分换成了搜索：在一族容易处理的分布里找离后验最近的那个，最大化 ELBO 等价于最小化反向 KL。这个方向的脾气会留下印记——它宁可缩进一个模态也不肯把质量摊到目标不支持的地方，于是不确定性被系统性地低估。这不是实现有 bug，是目标函数选了这个方向就会这样。

两条路上还有一条共同的诊断线索：重要性权重的方差。提议分布的尾部比目标薄时，权重按指数发散，表面十万个样本的有效样本量可能只有十几个。它在第一节以方差缩减的形式出现，在粒子滤波里以每步权重塌缩的形式出现，在检验变分近似质量时又被反过来用作探针。

六节的分工是这样的。第一节 \hyperref[sec:13-Machine-Learning/08-Sampling-and-Inference/01]{``Monte Carlo 用随机样本逼近期望''} 建立误差的记账方式，包括重要性采样的权重诊断。第二节 \hyperref[sec:13-Machine-Learning/08-Sampling-and-Inference/02]{``MCMC 在无法独立采样时构造平稳链''} 用细致平衡把归一化常数从接受率里约掉，并把遍历性、谱间隙、跨模态失败讲清楚。第三节 \hyperref[sec:13-Machine-Learning/08-Sampling-and-Inference/03]{``变分推断把后验近似变成优化''} 换到另一条路，把代价定位到变分族与 KL 方向上。第四节 \hyperref[sec:13-Machine-Learning/08-Sampling-and-Inference/04]{``Bootstrap 重建分析流程的不确定性''} 处理另一类问题：没有后验可谈，只问换一批数据整条流程会波动多少。第五节 \hyperref[sec:13-Machine-Learning/08-Sampling-and-Inference/05]{``粒子滤波用带权样本追踪非线性状态''} 把重要性采样接进时间递推。第六节 \hyperref[sec:13-Machine-Learning/08-Sampling-and-Inference/06]{``配分函数把能量变成全局概率''} 正面对付那个一直被约掉的常数，并说明为什么局部能量便宜而全局学习仍然依赖采样。

\subsection{读之前先把四类误差分开}

Monte Carlo 的有限样本误差、MCMC 的混合与自相关、变分族造成的近似偏差、Bootstrap 对重采样单位的依赖，这四者的补救动作完全不同，混起来看就会用错药：看到不确定性偏小去加迭代次数，看到区间偏窄去加 bootstrap 次数，多半都是在治错的那一层。粒子滤波与能量模型是前面工具的组合应用，可以第二遍再进。

读任何一份用了这些方法的报告时，问四个问题：目标分布是什么，样本是不是真的来自它，近似族排除了哪些形状，报告的区间描述的是统计不确定性还是数值计算误差。还有一条贯穿全章的提醒——ELBO 上升、链看着平稳、重构图像变好，这三样都不是正确性的证据，它们只是没报警而已。这一章所有的诊断工具都只能证伪，把这一点接受下来，后面六节读起来会顺很多。
